\documentclass[aspectratio=169,11pt]{beamer}

\usepackage[T1]{fontenc}
\usepackage[utf8]{inputenc}
\usepackage{lmodern}
\usepackage{microtype}
\usepackage{amsmath,amssymb,mathtools,bm}
\usepackage{booktabs,array}
\usepackage{xcolor}

\definecolor{TitleBlue}{HTML}{163A5F}
\definecolor{AccentTeal}{HTML}{2D6F73}
\definecolor{DeepNavy}{HTML}{102A43}
\definecolor{SoftBlue}{HTML}{EAF3F8}
\definecolor{SoftTeal}{HTML}{EDF8F6}
\definecolor{SoftGray}{HTML}{F4F6F8}
\definecolor{BodyGray}{HTML}{243B53}
\definecolor{WarningRed}{HTML}{9B2C2C}

\setbeamercolor{normal text}{fg=BodyGray,bg=white}
\setbeamercolor{frametitle}{fg=TitleBlue,bg=white}
\setbeamercolor{structure}{fg=AccentTeal}
\setbeamercolor{block title}{fg=white,bg=TitleBlue}
\setbeamercolor{block body}{fg=BodyGray,bg=SoftBlue}
\setbeamercolor{block title alerted}{fg=white,bg=WarningRed}
\setbeamercolor{block body alerted}{fg=BodyGray,bg=SoftGray}
\setbeamercolor{block title example}{fg=white,bg=AccentTeal}
\setbeamercolor{block body example}{fg=BodyGray,bg=SoftTeal}

\setbeamerfont{normal text}{size=\fontsize{15.5}{18.8}\selectfont}
% This is a dense personal-reference deck: frame titles deliberately use
% approximately the same size as the body, distinguished only by weight/color.
\setbeamerfont{frametitle}{size=\fontsize{16}{19}\selectfont,series=\bfseries}
\setbeamerfont{block title}{size=\fontsize{15.5}{18.5}\selectfont,series=\bfseries}
\setbeamerfont{block body}{size=\fontsize{15}{18}\selectfont}
\setbeamerfont{footline}{size=\fontsize{10}{12}\selectfont}

\setbeamertemplate{navigation symbols}{}
\setbeamertemplate{itemize item}{\color{AccentTeal}$\blacktriangleright$}
\setbeamertemplate{itemize subitem}{\color{AccentTeal}$\bullet$}
\setlength{\leftmargini}{1.25em}
\setlength{\leftmarginii}{1.15em}
\setlength{\parskip}{0.05em}

\newcommand{\bracket}[2]{\left[#1,#2\right]}
\newcommand{\dd}{\,\mathrm{d}}
\newcommand{\average}[1]{\left\langle #1\right\rangle}
\newcommand{\Rpsi}{R_{\psi}}
\newcommand{\Romega}{R_{\omega}}
\newcommand{\vect}[1]{\bm{#1}}
\newcommand{\currentmodule}{Module 1}
\newcommand{\moduleprogress}{1 of 6}
\newcommand{\setmodule}[2]{%
  \renewcommand{\currentmodule}{#1}%
  \renewcommand{\moduleprogress}{#2}%
}

\setbeamertemplate{frametitle}{%
  \vspace{0.02cm}
  \usebeamerfont{frametitle}\insertframetitle\par
  \vspace{-0.18cm}
  {\color{AccentTeal}\rule{\textwidth}{0.8pt}}
}

\setbeamertemplate{footline}{%
  \leavevmode
  \hbox{%
    \begin{beamercolorbox}[wd=.42\paperwidth,ht=3.1ex,dp=1.35ex,leftskip=.55cm]{author in head/foot}%
      \usebeamerfont{footline}\color{TitleBlue}\textbf{PINN\_M3DC1}\hspace{0.6em}|\hspace{0.6em}\currentmodule
    \end{beamercolorbox}%
    \begin{beamercolorbox}[wd=.58\paperwidth,ht=3.1ex,dp=1.35ex,rightskip=.55cm plus1fil]{date in head/foot}%
      \usebeamerfont{footline}\color{BodyGray}\moduleprogress\hspace{1.0em}\textbullet\hspace{1.0em}Slide \insertframenumber\ of \inserttotalframenumber
    \end{beamercolorbox}%
  }
}

\setbeamersize{text margin left=0.70cm,text margin right=0.70cm}
\AtBeginDocument{%
  \setlength{\abovedisplayskip}{0.30em}%
  \setlength{\belowdisplayskip}{0.30em}%
  \setlength{\abovedisplayshortskip}{0.20em}%
  \setlength{\belowdisplayshortskip}{0.20em}%
}
\title{Modules 1, 4, and 5 Summary}
\subtitle{Reduced resistive MHD, data-only surrogate, and physics-informed surrogate}
\author{PINN\_M3DC1 Project}
\date{4 August 2026}

\begin{document}

% --------------------------------------------------------------------------
% Module 1: six slides
% --------------------------------------------------------------------------

\setmodule{Module 1}{1 of 6}
\begin{frame}{Module 1: Reduced resistive MHD}
The version-1 model is two-dimensional, incompressible, two-field reduced resistive magnetohydrodynamics (RMHD) on a doubly periodic Cartesian domain:
\[
(x,y,t)\in[0,2\pi)\times[0,2\pi)\times[0,T].
\]
\[
\vect{q}(x,y,t)=\bigl(\psi,U\bigr),
\qquad
\vect{\mu}=\bigl(\eta,\nu,A_0\bigr).
\]

\begin{itemize}
  \item $\psi$ is the magnetic-flux function; $U$ is the flow stream function.  The coefficients $\eta$ and $\nu$ are constant within one case; $A_0$ indexes its initial perturbation.
  \item Frozen assumptions: $\partial_z=0$, constant density, $\nabla\!\cdot\vect v=0$, scalar resistivity and viscosity, and a prescribed uniform guide field outside the two evolved fields.
  \item Omitted physics includes compressibility, density/temperature evolution, Hall and two-fluid terms, kinetic closures, toroidal geometry, conducting walls, and free-boundary effects.
\end{itemize}
\textcolor{WarningRed}{\textbf{Scope:}} these are synthetic RMHD trajectories in an M3D-C1-shaped workflow, not genuine M3D-C1 output.
\end{frame}

\setmodule{Module 1}{2 of 6}
\begin{frame}{Two potentials define the physical fields}
\begin{columns}[T,onlytextwidth]
\begin{column}{0.49\textwidth}
\textbf{Magnetic field}
\[
\vect{B}_{\perp}
=\nabla\psi\times\widehat{\vect{z}}
=\bigl(\psi_y,-\psi_x\bigr),
\]
\[
J=-\nabla^2\psi.
\]
$J\widehat{\vect z}$ is the normalized out-of-plane current variable; a uniform guide field $B_g\widehat{\vect z}$ is represented separately.
\end{column}
\begin{column}{0.49\textwidth}
\textbf{Fluid velocity}
\[
\vect{v}_{\perp}
=\widehat{\vect{z}}\times\nabla U
=\bigl(-U_y,U_x\bigr),
\]
\[
\omega=\nabla^2U.
\]
$\omega\widehat{\vect z}=\nabla\times\vect v_\perp$ is the out-of-plane vorticity.  Adding a spatial constant to $U$ changes no velocity.
\end{column}
\end{columns}

\[
\nabla\cdot\vect{B}_{\perp}=0,
\qquad
\nabla\cdot\vect{v}_{\perp}=0,
\qquad
\average{J}=0.
\]
Thus both divergence constraints hold identically. The network predicts $(\psi,U)$; derivatives supply $(\vect B_\perp,\vect v_\perp,J,\omega)$.  Recovering $U$ from $\omega=\nabla^2U$ requires fixing its zero Fourier mode.
\end{frame}

\setmodule{Module 1}{3 of 6}
\begin{frame}{The Poisson bracket represents advection}
\[
\boxed{\bracket{f}{g}=f_xg_y-f_yg_x.}
\]

The stream-function velocity converts transport into a bracket:
\[
\vect{v}_{\perp}\cdot\nabla f
=(-U_y)f_x+(U_x)f_y
=\bracket{U}{f},
\qquad
\frac{Df}{Dt}=f_t+\bracket{U}{f}.
\]

\[
\bracket{f}{g}=-\bracket{g}{f},
\qquad
\int_{\Omega}\bracket{f}{g}\dd A=0,
\qquad
\int_\Omega h\bracket{f}{g}\dd A
=\int_\Omega f\bracket{g}{h}\dd A.
\]
The integral identities follow from periodicity.  They make the nonlinear terms conservative: brackets redistribute structure and exchange magnetic and kinetic energy, but create no net dissipation.

In the ideal flux equation ($\eta=0$), $D\psi/Dt=0$: magnetic-flux labels move with the fluid.  Reconnection therefore requires a non-ideal term or a breakdown of the model assumptions.
\end{frame}

\setmodule{Module 1}{4 of 6}
\begin{frame}{Two PDEs govern the dynamics}
\begin{align*}
\boxed{\psi_t+\bracket{U}{\psi}=\eta\nabla^2\psi,}
\qquad
\boxed{\omega_t+\bracket{U}{\omega}=\bracket{J}{\psi}+\nu\nabla^2\omega.}
\end{align*}

\begin{itemize}
  \item $\bracket{U}{\psi}$ transports magnetic-flux contours with the fluid.
  \item $\eta\nabla^2\psi=-\eta J$ permits field-line slippage.  In a sheet of width $\delta$, $\eta\nabla^2\psi\sim\eta\psi/\delta^2$; small global $\eta$ can therefore matter locally.
  \item $\bracket{U}{\omega}$ transports vorticity; $\bracket{J}{\psi}=J_x\psi_y-J_y\psi_x$ is the curl of the in-plane Lorentz force.
  \item $\nu\nabla^2\omega$ damps fine-scale flow. Pressure is not zero; curling the constant-density momentum equation eliminates its gradient.
\end{itemize}
Evolving $(\psi,\omega)$ requires solving $\nabla^2U=\omega$ each step; the neural surrogate outputs $(\psi,U)$ directly.
\end{frame}

\setmodule{Module 1}{5 of 6}
\begin{frame}{Alfv\'en normalization sets the scales}
\[
V_A=\frac{B_0}{\sqrt{\mu_0\rho_0}},
\qquad
t_A=\frac{L_0}{V_A}.
\]

\[
\psi^*=\frac{\psi_{\rm dim}}{B_0L_0},
\qquad
U^*=\frac{U_{\rm dim}}{L_0V_A},
\qquad
t^*=\frac{t_{\rm dim}}{t_A}.
\]

\[
\eta=\frac{\eta_m}{L_0V_A},
\qquad
\nu=\frac{\nu_d}{L_0V_A}.
\]

\[
S=\frac{1}{\eta}=\frac{t_\eta}{t_A},\quad t_\eta=\frac{L_0^2}{\eta_m},
\qquad
\mathrm{Re}=\frac{1}{\nu}=\frac{t_\nu}{t_A},\quad t_\nu=\frac{L_0^2}{\nu_d},
\qquad
\mathrm{Pm}=\frac{\nu}{\eta}.
\]

$S$ and $\mathrm{Re}$ compare diffusion times with the Alfv\'en time; $\mathrm{Pm}$ compares momentum and magnetic diffusion.  The identities $S=1/\eta$ and $\mathrm{Re}=1/\nu$ hold under this declared scaling.  Neural standardization is a later numerical transform, not a change of physical units.
\end{frame}

\setmodule{Module 1}{6 of 6}
\begin{frame}{Periodicity and energy close the model}
For $f\in\{\psi,U\}$, periodic values and every PDE-required derivative match across opposite faces. Gauge and mean-flux conditions are
\[
\boxed{\average{U}(t)=0,}
\qquad
\boxed{\frac{\mathrm d}{\mathrm dt}\average{\psi}=0.}
\]
The exact periodic energy law is
\[
\boxed{
\frac{\mathrm d}{\mathrm dt}\frac12\int_{\Omega}
\left(|\nabla\psi|^2+|\nabla U|^2\right)\dd A
=-\eta\int_{\Omega}J^2\dd A
-\nu\int_{\Omega}\omega^2\dd A\le 0.}
\]
The left side contains magnetic plus kinetic energy; only resistivity and viscosity cause net decay.  This identity is a solver/PINN verification test, but it is not sufficient by itself: spatially wrong fields can share the same integral energy.

Island coalescence exercises nonlinear advection, Lorentz forcing, current-sheet formation, resistive topology change, and sensitivity to $(\eta,\nu,A_0)$.  Its exact equilibrium, perturbation, and reconnection diagnostic are Module~2 contracts.
\end{frame}

% --------------------------------------------------------------------------
% Module 4: four slides
% --------------------------------------------------------------------------

\setmodule{Module 4}{1 of 4}
\begin{frame}{Module 4: A data-only field surrogate}
The parameterized coordinate network approximates a family of fields:
\[
\boxed{
\mathcal N_{\vect\theta}:
(x,y,t,\eta,\nu,A_0)
\longmapsto
\bigl(\widehat\psi,\widehat U\bigr).}
\]

Complete-case split: $\vect\mu_c=(\eta_c,\nu_c,A_{0,c})$.
\begin{itemize}
  \item No point from a held-out case enters training or preprocessing.
  \item A coordinate network evaluates any requested coordinate directly; it is not autoregressive and does not propagate its previous prediction.
  \item It is not a general neural operator: the initial condition varies only through scalar $A_0$, not an arbitrary function.
  \item Parameter interpolation and extrapolation are separate; evaluation inside $t\in[0,T]$ is not forecasting beyond $T$. Many points from few cases can still mean memorization.
\end{itemize}
\end{frame}

\setmodule{Module 4}{2 of 4}
\begin{frame}{Periodic inputs respect the domain}
\[
\phi_x=\frac{2\pi x}{L_x},
\qquad
\phi_y=\frac{2\pi y}{L_y},
\]
\[
\vect z=
\bigl(
\sin\phi_x,\cos\phi_x,
\sin\phi_y,\cos\phi_y,
\widetilde t,\widetilde r_\eta,\widetilde r_\nu,\widetilde A_0
\bigr).
\]

\[
r_\eta=\log_{10}\eta,
\qquad
r_\nu=\log_{10}\nu,
\qquad
\widetilde a=\frac{a-\mu_a^{\rm train}}{\sigma_a^{\rm train}}.
\]

Smooth functions of $\vect z$ are periodic in value and every derivative.  Scaling statistics and output standardization are computed from training cases only, then frozen for validation, testing, and deployment.

\textbf{Baseline:} six fully connected $\tanh$ layers of width $128$ and a linear two-output head.  The smooth activation also supports the high derivatives required later.  Module~5 retains the same architecture family and periodic representation so physics information is the controlled change.
\end{frame}

\setmodule{Module 4}{3 of 4}
\begin{frame}{Training uses labeled field error only}
Each supervised record retains case identity:
\[
\bigl(x_n,y_n,t_n,\vect\mu_{c(n)};\ \psi_n,U_n\bigr),
\qquad c(n)=\text{source case}.
\]
For standardized output errors
\[
e_{\psi,n}=\widehat{\widetilde\psi}_n-\widetilde\psi_n,
\qquad
e_{U,n}=\widehat{\widetilde U}_n-\widetilde U_n,
\]
the case-balanced objective is
\[
\boxed{
\mathcal L_{\mathrm{train}}=\mathcal L_{\mathrm{data}}
=\frac{1}{|\mathcal C_{\mathcal B}|}
\sum_{c\in\mathcal C_{\mathcal B}}
\left[
\frac{1}{N_c}\sum_{n\in c}e_{\psi,n}^2
+\frac{1}{N_c}\sum_{n\in c}e_{U,n}^2
\right].}
\]

Case balancing prevents long trajectories from dominating. No physics-derived constraint or target appears: no PDE, IC, BC, gauge, global-balance, $J$, $\omega$, derivative, or spectral loss. Optimizer and regularization choices do not make the model physics-informed.
\end{frame}

\setmodule{Module 4}{4 of 4}
\begin{frame}{The baseline makes the PINN testable}
Evaluation uses complete unseen fields, not training loss alone. For $a\in\{\psi,U,J,\omega\}$,
\[
\varepsilon_{L_2}^{(a)}
=\frac{\|\widehat a-a\|_2}{\|a\|_2+\epsilon}.
\]

\[
\widehat J=-\nabla^2\widehat\psi,
\qquad
\widehat\omega=\nabla^2\widehat U,
\]
and current-sheet sensitivity is measured by
\[
J_{\max}(t)=\max_{(x,y)\in\Omega}|J(x,y,t)|.
\]
Derived fields, spectra, energy, reconnection measures, periodicity defects, and term-normalized PDE residuals are evaluation-only diagnostics; none are backpropagated.

\textcolor{WarningRed}{\textbf{Controlled comparison:}} Modules~4 and~5 share case splits, labeled-data budgets, preprocessing, periodic encoding, architecture family, data loss, metrics, and random seeds.  Report interpolation, extrapolation, data-efficiency curves, repeated seeds, and wall-clock/compute cost. Physics helps only if a matched PINN improves a declared metric.
\end{frame}

% --------------------------------------------------------------------------
% Module 5: five slides
% --------------------------------------------------------------------------

\setmodule{Module 5}{1 of 5}
\begin{frame}{Module 5: PDE residuals enter training}
The network still predicts $(\widehat\psi,\widehat U)$, while automatic differentiation constructs
\[
\widehat J=-\nabla^2\widehat\psi,
\qquad
\widehat\omega=\nabla^2\widehat U.
\]

The local reduced-MHD equation defects are
\[
\boxed{
\Rpsi
=\widehat\psi_t
+\bracket{\widehat U}{\widehat\psi}
-\eta\nabla^2\widehat\psi,}
\]
\[
\boxed{
\Romega
=\widehat\omega_t
+\bracket{\widehat U}{\widehat\omega}
-\bracket{\widehat J}{\widehat\psi}
-\nu\nabla^2\widehat\omega.}
\]

For example, $\bracket{\widehat U}{\widehat\psi}=\widehat U_x\widehat\psi_y-\widehat U_y\widehat\psi_x$.  Automatic differentiation (AD) differentiates the implemented neural computational graph, not the unknown physical solution.

At unlabeled collocation points $(x_f,y_f,t_f,\vect\mu_f)$, the PINN receives coordinates and coefficients but no solver field values.  The learning signal is the defect of the declared equation.  A small sampled residual is evidence only on the sampled residual norm; it is not automatically a small field error.
\end{frame}

\setmodule{Module 5}{2 of 5}
\begin{frame}{The objective combines five constraints}
\[
\boxed{
\mathcal L_{\mathrm{total}}
=w_{\mathrm{data}}\mathcal L_{\mathrm{data}}
+w_{\mathrm{PDE}}\mathcal L_{\mathrm{PDE}}
+w_{\mathrm{IC}}\mathcal L_{\mathrm{IC}}
+w_{\mathrm{BC}}\mathcal L_{\mathrm{BC}}
+w_{\mathrm{global}}\mathcal L_{\mathrm{global}}.}
\]

\[
\mathcal L_{\mathrm{PDE}}
=\frac{1}{N_f}\sum_{i=1}^{N_f}
\left[
\left(\frac{\Rpsi^{(i)}}{s_\psi}\right)^2
+\left(\frac{\Romega^{(i)}}{s_\omega}\right)^2
\right].
\]

\begin{itemize}
  \item $\mathcal L_{\mathrm{data}}$: the identical labeled observations used by Module 4.
  \item $\mathcal L_{\mathrm{IC}}$ at $t=0$ selects the trajectory; the PDE alone is not unique.
  \item Hard sine-cosine periodicity gives $\mathcal L_{\mathrm{BC}}=0$; $\mathcal L_{\mathrm{global}}$ contains gauge, mean-flux, and optional balances.
\end{itemize}
Data, collocation, IC, boundary, and quadrature points are distinct populations. Normalize complete residual equations, not individual terms whose cancellation carries the physics.
\end{frame}

\setmodule{Module 5}{3 of 5}
\begin{frame}{The vorticity residual reaches fourth order}
Because $\widehat\omega=\nabla^2\widehat U$,
\[
\widehat\omega_t\sim\partial_t\partial_x^2\widehat U,
\qquad
\bracket{\widehat U}{\widehat\omega}\sim\partial_x^3\widehat U,
\qquad
\nabla^2\widehat\omega=\nabla^4\widehat U.
\]
Third derivatives of $\widehat\psi$ also occur.  Smooth activations such as $\tanh$, sufficient precision, and analytic derivative tests are mandatory; ReLU is unsuitable for this strong-form residual.

For $q=s_q\widetilde q+\mu_q$ and $\widetilde x=(x-\mu_x)/s_x$,
\[
q_x=\frac{s_q}{s_x}\widetilde q_{\widetilde x},
\qquad
q_{xxxx}=\frac{s_q}{s_x^4}\widetilde q_{\widetilde x\widetilde x\widetilde x\widetilde x}.
\]
Inverse output scaling and coordinate chain rules must stay inside the residual graph; omitting them silently changes the PDE.

The PDE cannot identify a spatially constant offset in $U$, so the independent gauge condition $\boxed{\average{\widehat U}(t,\vect\mu)=0}$ remains necessary.
\end{frame}

\setmodule{Module 5}{4 of 5}
\begin{frame}{Balances and curricula stabilize training}
Periodic global checks include $\frac{\mathrm d}{\mathrm dt}\average{\widehat\psi}=0$ and $\average{\widehat U}=0$.
\[
\boxed{
\dot E
+\eta\int_{\Omega}\widehat J^2\dd A
+\nu\int_{\Omega}\widehat\omega^2\dd A=0,
\qquad
E=\frac12\int_\Omega\left(|\nabla\widehat\psi|^2+|\nabla\widehat U|^2\right)\dd A.}
\]

Verified quadrature is required; global penalties supplement local residuals. Collocation defines the minimized norm: augment uniform samples adaptively, but keep validation points fixed.

\textbf{Curriculum:} derivative tests $\rightarrow$ data warm start $\rightarrow\Rpsi\rightarrow\Romega$
$\rightarrow$ time curriculum $\rightarrow$ parameterized cases.

Raw loss magnitudes do not diagnose balance. Monitor gradient norms and cosine similarities; opposing gradients can reveal scale conflict, finite capacity, noise, or model discrepancy.
\end{frame}

\setmodule{Module 5}{5 of 5}
\begin{frame}{Physics must add measurable value}
Unseen-case scoring covers field errors, residuals, condition defects, global balances, and cost.

\begin{itemize}
  \item A PINN may improve labeled-data efficiency, derivative fidelity, parameter generalization, or physical admissibility.
  \item It may also train more slowly, oversmooth current sheets, suffer gradient conflict, or inherit error from an incomplete physical model.
  \item Reports include repeated seeds, failures, extrapolation, and data-versus-compute fairness.
\end{itemize}

\textbf{Minimum ablation ladder:}
\[
\mathcal L_{\rm data}
\;\longrightarrow\;
\mathcal L_{\rm data}+\mathcal L_{\rm PDE}
\;\longrightarrow\;
\mathcal L_{\rm data}+\mathcal L_{\rm PDE}
+\mathcal L_{\rm IC}+\mathcal L_{\rm global},
\]
Vary one component, fix the label budget, and disclose extra collocation/optimization cost.

\textcolor{WarningRed}{\textbf{Scientific conclusion:}} no win is assumed.  The result must be improved held-out accuracy, admissibility, or data efficiency---not merely lower training loss.  Scope remains synthetic two-field RMHD; validated M3D-C1 claims require genuine data and Module~7.
\end{frame}

\end{document}
